\chapter{Special Elements}
\label{chap:SpecialElements}
The Schema definition of the XML parameter file currently allows for two
special elements to be specified. These are
\begin{itemize}
\item \xel{expert}
\item \xel{devel}
\end{itemize}
These are first level elements coming at the end of the parameter file, see
Sect.~\ref{sec:general}.
\subsubsection*{\xel{expert}}
The \xel{expert} element is an empty element with one attribute called
\xat{override}. Specifying \xml{<expert override="yes">} in the parameter
file
turns off the Expert() module of the CFS++ --- OLAS interface. This expert
module normally tries to set certain
parameters to sensible values depending on the values of other parameters or
to fix inconsistencies in the input parameters. E.g.~the expert module will
turn off preconditioning, if a direct solver is specified. There may be situations, however, where this behaviour is not
desired, e.g.~one may want to test the performance of an ILU preconditioner
without the \xml{Metis} matrix re-ordering the expert would trigger by default.
Note that turning off the expert module may lead to run-time aborts in OLAS due
to inconsistent settings. So be sure to really be an expert, if you do this.
\subsubsection*{\xel{devel}}
The \xel{devel} element can have arbitrary content, but no children or
attributes. It is solely there for the sake of the developer who wants to
quickly test something, without the need to alter the Schema definition itself.
For normal users the element is not of interest, since it is never accessed
in the standard CFS++ source code.
